\section{Równoważności kategorii}\label{sec:equivalence}

W algebrze odróżniamy równość od izomorfizmu: przestrzeń $V$ może być izomorficzna z $k^n$, choć nie jest tym samym zbiorem. Podobne zjawisko występuje na poziomie całych kategorii. Równość kategorii jest zwykle zbyt rygorystyczna, dlatego wprowadzamy równoważność.

\begin{definicja}[Równoważność kategorii]\label{def:equivalence}
Funktor $F:\mathcal C\to\mathcal D$ jest \emph{równoważnością kategorii}, jeśli istnieje funktor $G:\mathcal D\to\mathcal C$ oraz izomorfizmy naturalne $G\circ F\cong\id_{\mathcal C}$ i $F\circ G\cong\id_{\mathcal D}$. Mówimy wtedy, że $\mathcal C$ i $\mathcal D$ są \emph{równoważne}, i piszemy $\mathcal C\simeq\mathcal D$. Funktor $G$ jest \emph{quasi-odwrotnością} $F$.
\end{definicja}

Nie wymagamy, aby $G\circ F$ oraz $F\circ G$ były \emph{równe} odpowiednim identycznościom. Wystarcza, że można je z identycznościami porównać przez izomorfizmy naturalne.

\begin{definicja}[Istotna surjektywność na obiektach]
Funktor $F:\mathcal C\to\mathcal D$ jest \emph{istotnie surjektywny na obiektach}, jeśli dla każdego obiektu $Y\in\Ob(\mathcal D)$ istnieje $X\in\Ob(\mathcal C)$ oraz izomorfizm $F(X)\cong Y$ w $\mathcal D$.
\end{definicja}

\begin{twierdzenie}[Kryterium równoważności]\label{tw:equivalence}
Niech $\mathcal C$ i $\mathcal D$ będą małymi kategoriami. Funktor $F:\mathcal C\to\mathcal D$ jest równoważnością wtedy i tylko wtedy, gdy jest w pełni wierny i istotnie surjektywny na obiektach.
\end{twierdzenie}
\begin{proof}
Załóżmy najpierw, że $F$ jest równoważnością. Wybierzmy quasi-odwrotność $G$ i izomorfizmy naturalne $\alpha:GF\Rightarrow\id_{\mathcal C}$ oraz $\beta:FG\Rightarrow\id_{\mathcal D}$. Dla każdego $Y\in\Ob(\mathcal D)$ morfizm $\beta_Y:FG(Y)\to Y$ jest izomorfizmem, co dowodzi istotnej surjektywności. Jeśli $f,g:A\to B$ i $F(f)=F(g)$, to $GF(f)=GF(g)$. Naturalność $\alpha$ daje $\alpha_BGF(f)=f\alpha_A$ oraz $\alpha_BGF(g)=g\alpha_A$; ponieważ $\alpha_A$ jest odwracalny, wynika $f=g$. To dowodzi wierności.

Sprawdźmy pełność bez nieuzasadnionego zakładania, że $F(\alpha_X)=\beta_{F(X)}$. Dla $X\in\Ob(\mathcal C)$ zdefiniujmy automorfizm $c_X:F(X)\to F(X)$ przez $c_X=F(\alpha_X)\circ\beta_{F(X)}^{-1}$. Niech $u:F(A)\to F(B)$ będzie dowolnym morfizmem. Połóżmy $u'=c_B^{-1}\circ u\circ c_A$ i $f=\alpha_B\circ G(u')\circ\alpha_A^{-1}:A\to B$. Naturalność $\beta$ daje
$FG(u')=\beta_{F(B)}^{-1}\circ u'\circ\beta_{F(A)}$, zatem
\[
F(f)=F(\alpha_B)\circ FG(u')\circ F(\alpha_A)^{-1}
=c_B\circ u'\circ c_A^{-1}=u.
\]
Każdy $u$ ma więc przeciwobraz przez $F_{A,B}$; funktor $F$ jest pełny.

Na odwrót, załóżmy pełną wierność i istotną surjektywność $F$. Dla każdego $Y\in\Ob(\mathcal D)$ wybierzmy obiekt $G(Y)\in\Ob(\mathcal C)$ i izomorfizm $\varepsilon_Y:F(G(Y))\to Y$. Dla każdego $Y$ zbiór dopuszczalnych par $(X,\varepsilon)$ jest niepusty dzięki istotnej surjektywności. Jest zbiorem, ponieważ obie kategorie są małe. Rodzina tych zbiorów jest indeksowana zbiorem $\Ob(\mathcal D)$; aksjomat wyboru daje jednoczesny wybór par $(G(Y),\varepsilon_Y)$. Dla strzałki $u:Y\to Y'$ pełność i wierność dają dokładnie jeden morfizm $G(u):G(Y)\to G(Y')$ spełniający
\[
F(G(u))=\varepsilon_{Y'}^{-1}\circ u\circ\varepsilon_Y.
\]
Dla identyczności mamy $F(G(\id_Y))=\varepsilon_Y^{-1}\circ\varepsilon_Y=\id_{F(G(Y))}=F(\id_{G(Y)})$, więc wierność daje $G(\id_Y)=\id_{G(Y)}$. Jeśli $u:Y\to Y'$ i $v:Y'\to Y''$, to
\[
\begin{aligned}
F(G(v)\circ G(u))
&=(\varepsilon_{Y''}^{-1}\circ v\circ\varepsilon_{Y'})
  \circ(\varepsilon_{Y'}^{-1}\circ u\circ\varepsilon_Y)\\
&=\varepsilon_{Y''}^{-1}\circ(v\circ u)\circ\varepsilon_Y\\
&=F(G(v\circ u)).
\end{aligned}
\]
Obie strzałki pod funktorem $F$ należą do $\Hom(G(Y),G(Y''))$, więc z wierności $G(v\circ u)=G(v)\circ G(u)$. Tak powstaje funktor $G:\mathcal D\to\mathcal C$. Mnożąc definiującą równość z lewej przez $\varepsilon_{Y'}$, otrzymujemy $\varepsilon_{Y'}\circ FG(u)=u\circ\varepsilon_Y$. Jest to naturalność $\varepsilon:FG\Rightarrow\id_{\mathcal D}$, a każda składowa jest izomorfizmem z wyboru.

Dla $A\in\Ob(\mathcal C)$ pełna wierność wyznacza dokładnie jeden $\eta_A:A\to GF(A)$ z $F(\eta_A)=\varepsilon_{F(A)}^{-1}$. Z pełności istnieje także $\theta_A:GF(A)\to A$ z $F(\theta_A)=\varepsilon_{F(A)}$. Zatem
\[
\begin{aligned}
F(\theta_A\circ\eta_A)&=\varepsilon_{F(A)}\circ\varepsilon_{F(A)}^{-1}=F(\id_A),\\
F(\eta_A\circ\theta_A)&=\varepsilon_{F(A)}^{-1}\circ\varepsilon_{F(A)}=F(\id_{GF(A)}).
\end{aligned}
\]
Wierność daje $\theta_A\circ\eta_A=\id_A$ i $\eta_A\circ\theta_A=\id_{GF(A)}$. Tym samym $\eta_A$ jest izomorfizmem o odwrotności $\theta_A$. Dla $f:A\to A'$ naturalność $\varepsilon$ daje
$FGF(f)\circ\varepsilon_{F(A)}^{-1}=\varepsilon_{F(A')}^{-1}\circ F(f)$. Stąd
\[
\begin{aligned}
F(GF(f)\circ\eta_A)
&=FGF(f)\circ\varepsilon_{F(A)}^{-1}\\
&=\varepsilon_{F(A')}^{-1}\circ F(f)\\
&=F(\eta_{A'}\circ f).
\end{aligned}
\]
Z wierności $GF(f)\circ\eta_A=\eta_{A'}\circ f$. Zatem $\eta:\id_{\mathcal C}\Rightarrow GF$ jest izomorfizmem naturalnym, a $F$ ma quasi-odwrotność.
\end{proof}


\begin{uwaga}[Zakres założenia o małości]
Małość kategorii w twierdzeniu~\ref{tw:equivalence} jest wygodnym warunkiem uzasadniającym jednoczesny wybór reprezentantów. W praktyce kryterium stosuje się również do kategorii lokalnie małych, przy odpowiedniej konwencji dotyczącej wyborów dla dużych klas. Kategoria wszystkich skończenie wymiarowych przestrzeni wektorowych nie jest dosłownie mała; można jednak pracować z jej małą podkategorią reprezentantów $k^0,k^1,k^2,\ldots$ lub przeprowadzić konstrukcję w ustalonym uniwersum zbiorów. W przykładzie macierzy rozumiemy zatem równoważność w ustalonym uniwersum (z wyborem w większym uniwersum) albo jako równoważność z małym modelem. Bez konwencji umożliwiającej wybór klasowy nie wnioskujemy automatycznie o istnieniu globalnego quasi-odwrotnego funktora na właściwej klasie obiektów.
\end{uwaga}


\begin{przyklad}[Macierze i przestrzenie skończenie wymiarowe]\label{przyklad:mat-equivalence}
Funktor $F:\cat{Mat}_k\to\Vect_k^{\mathrm{fd}}$ z przykładu~\ref{przyklad:mat} jest w pełni wierny, bo każdemu odwzorowaniu liniowemu $k^n\to k^m$ odpowiada dokładnie jedna macierz. Jest też istotnie surjektywny, ponieważ każda skończenie wymiarowa przestrzeń $V$ jest izomorficzna z $k^{\dim V}$. Zatem
\[
\cat{Mat}_k\simeq\Vect_k^{\mathrm{fd}}.
\]
Kategorie te nie są jednak dosłownie równe: w pierwszej obiektami są liczby naturalne, a w drugiej przestrzenie wektorowe (zwykle rozpatrywane bez uprzedniego wyboru baz). To modelowa sytuacja, w której równoważność jest właściwszym pojęciem niż równość.
\end{przyklad}

\begin{przyklad}[Szkielet kategorii]
\emph{Szkieletem} kategorii $\mathcal C$ nazywamy pełną podkategorię zawierającą dokładnie po jednym reprezentancie każdej klasy izomorfizmu obiektów (o ile taki wybór można przeprowadzić). Dla $\Vect_k^{\mathrm{fd}}$ można wybrać reprezentantów $k^0,k^1,k^2,\ldots$, co daje kategorię równoważną $\cat{Mat}_k$. Każda mała kategoria jest równoważna swojemu szkieletowi, lecz szkielet nie musi być dosłownie jednoznaczny. W przykładzie przestrzeni wektorowych, dla pełnej formalności teoriomnogościowej, można pracować w ustalonym uniwersum zbiorów albo zastąpić rozważaną kategorię małym modelem zawierającym reprezentantów wszystkich jej klas izomorfizmu.
\end{przyklad}

\begin{stwierdzenie}[Szkielet małej kategorii]\label{cat:prop:skeleton}
Przy aksjomacie wyboru każda mała kategoria ma szkielet. Inkluzja szkieletu jest równoważnością, a dowolne dwa szkielety tej samej małej kategorii są izomorficzne jako kategorie.
\end{stwierdzenie}
\begin{proof}
Relacja izomorfizmu na zbiorze obiektów jest równoważnością: identyczność daje zwrotność, odwrotność izomorfizmu symetrię, a złożenie izomorfizmów przechodniość. Aksjomat wyboru pozwala wybrać jeden obiekt z każdej niepustej klasy. Pełna podkategoria $\mathcal S$ na wybranych obiektach jest szkieletem. Jej inkluzja jest w pełni wierna z definicji pełnej podkategorii oraz istotnie surjektywna dzięki wyborowi reprezentantów. Stosujemy twierdzenie~\ref{tw:equivalence}.

Dla dwóch szkieletów $\mathcal S,\mathcal T$ tej samej kategorii złożenie inkluzji $\mathcal S\to\mathcal C$ z quasi-odwrotnością inkluzji $\mathcal T\to\mathcal C$ daje funktor $H:\mathcal S\to\mathcal T$ w pełni wierny i istotnie surjektywny: złożenie bijekcji na zbiorach morfizmów jest bijekcją, a istotna surjektywność złożenia wynika z tego, że funktory zachowują izomorfizmy. Quasi-odwrotność inkluzji jest sama równoważnością na mocy definicji~\ref{def:equivalence}.

W kategorii $\mathcal T$ obiekty izomorficzne są równe, więc istotna surjektywność $H$ daje rzeczywistą surjektywność na obiektach. Jeżeli $H(S)=H(S')$, to pełna wierność odzwierciedla izomorfizmy: podnosimy identyczność w obu kierunkach, a ich złożenia są identycznościami dzięki wierności (por. stwierdzenie~\ref{tw:ff-reflect-iso}). Stąd $S\cong S'$, a ponieważ $\mathcal S$ jest szkieletem, mamy $S=S'$. Funktor $H$ jest więc bijekcją na obiektach i na każdym zbiorze morfizmów. Definiujemy jego ścisłą odwrotność przez odwrotność bijekcji na obiektach i jedyne przeciwobrazy morfizmów; zachowanie identyczności i złożeń wynika z ich zachowania przez $H$ oraz jednoznaczności przeciwobrazów. Oba złożenia są identycznościami na obiektach i morfizmach, czyli $H$ jest izomorfizmem kategorii.
\end{proof}

\begin{uwaga}[Kiedy używać jakiego pojęcia?]
Izomorfizm $A\cong B$ porównuje \emph{obiekty tej samej kategorii}. Izomorfizm naturalny $F\cong G$ porównuje \emph{funktory z tymi samymi dziedziną i przeciwdziedziną}. Równoważność $\mathcal C\simeq\mathcal D$ porównuje \emph{całe kategorie}. Są to trzy różne poziomy pojęcia ``to samo z dokładnością do odpowiedniej struktury''.
\end{uwaga}

